%%%%%%%%%%%%%%%%%%%%%%%%%%%%%%%%%%%%%%%%%%%%%%%%%%%%%%%%%%%%%%%%%%%%%%%%
% Some of my favourite personal adjustments
%%%%%%%%%%%%%%%%%%%%%%%%%%%%%%%%%%%%%%%%%%%%%%%%%%%%%%%%%%%%%%%%%%%%%%%%
%
% These are the adjustments that I consider necessary for typesetting
% a nice thesis. However, they are *not* included in the template, as
% I do not want to force you to use them.

% This ensures that I am able to typeset bold font in table while still aligning the numbers
% correctly.
\usepackage{etoolbox}
\usepackage{enumerate}

\usepackage[binary-units=true]{siunitx}
\DeclareSIUnit\px{px}

\sisetup{%
  detect-all           = true,
  detect-family        = true,
  detect-mode          = true,
  detect-shape         = true,
  detect-weight        = true,
  detect-inline-weight = math,
}

\newcommand{\overbar}[1]{\mkern 1.5mu\overline{\mkern-1.5mu#1\mkern-1.5mu}\mkern 1.5mu}

%%%%%%%%%%%%%%%%%%%%%%%%%%%%%%%%%%%%%%%%%%%%%%%%%%%%%%%%%%%%%%%%%%%%%%%%
% Hyperlinks & bookmarks
%%%%%%%%%%%%%%%%%%%%%%%%%%%%%%%%%%%%%%%%%%%%%%%%%%%%%%%%%%%%%%%%%%%%%%%%

\usepackage[%
  colorlinks = true,
  citecolor  = black,
  linkcolor  = RoyalBlue,
  urlcolor   = black,
  backref    = page
  ]{hyperref}

\usepackage[capitalise]{cleveref}

\renewcommand*{\backref}[1]{}
\renewcommand*{\backrefalt}[4]{%
    \ifcase #1%
          \or [cited on page~#2]%
          \else [cited on pages~#2]%
    \fi%
    }

%%%%%%%%%%%%%%%%%%%%%%%%%%%%%%%%%%%%%%%%%%%%%%%%%%%%%%%%%%%%%%%%%%%%%%%%
% Bibliography
%%%%%%%%%%%%%%%%%%%%%%%%%%%%%%%%%%%%%%%%%%%%%%%%%%%%%%%%%%%%%%%%%%%%%%%%

\usepackage[sort]{natbib}
\setcitestyle{%
    authoryear,%
    open={[},close={]},citesep={;},%
    aysep={},yysep={,},%
    notesep={, }}

% Explain list of backreferences.
% https://tex.stackexchange.com/a/70953/1340
\renewcommand{\bibpreamble}{%
  [Citing pages are listed after each reference.]%
  \par\bigskip}

% ACM bibliography style looks more professional
\bibliographystyle{ACM-Reference-Format}


%%%%%%%%%%%%%%%%%%%%%%%%%%%%%%%%%%%%%%%%%%%%%%%%%%%%%%%%%%%%%%%%%%%%%%%%
% Fonts
%%%%%%%%%%%%%%%%%%%%%%%%%%%%%%%%%%%%%%%%%%%%%%%%%%%%%%%%%%%%%%%%%%%%%%%%

\ifxetexorluatex
  \setmainfont{Minion Pro}
\else
  \usepackage[lf]{ebgaramond}
  \usepackage[oldstyle,scale=0.7]{sourcecodepro}
  \singlespacing
\fi

\renewcommand{\th}{\textsuperscript{\textup{th}}\xspace}

\usepackage{xeCJK}
% \setCJKmainfont{喜鹊聚珍体}

% \newacronym[description={Principal component analysis}]{PCA}{PCA}{principal component analysis}
% \newacronym                                            {SNF}{SNF}{Smith normal form}
% \newacronym[description={Topological data analysis}]   {TDA}{TDA}{topological data analysis}

% \newglossaryentry{LaTeX}{%
%   name        = {\LaTeX},
%   description = {A document preparation system},
%   sort        = {LaTeX},
% }

% \newglossaryentry{Real numbers}{%
%   name        = {$\real$},
%   description = {The set of real numbers},
%   sort        = {Real numbers},
% }

% \makeindex
% \makeglossaries

%%%%%%%%%%%%%%%%%%%%%%%%%%%%%%%%%%%%%%%%%%%%%%%%%%%%%%%%%%%%%%%%%%%%%%%%
% Load up my personal packages
%%%%%%%%%%%%%%%%%%%%%%%%%%%%%%%%%%%%%%%%%%%%%%%%%%%%%%%%%%%%%%%%%%%%%%%%

\usepackage{metalogo}


% Basics
\usepackage{fixltx2e}
\usepackage{url}
\usepackage{fancyvrb}
\usepackage{mdwlist}  % Miscellaneous list-related commands
\usepackage{bookmark}

% https://www.nesono.com/?q=book/export/html/347
% Package for inserting TODO statements in nice colorful boxes - so that you
% won’t forget to fix/remove them. To add a todo statement, use something like
% \todo{Find better wording here}.
\usepackage{todonotes}

%% Math
\usepackage{bm}       % Bold symbols in maths mode
\usepackage{amssymb}

% http://tex.stackexchange.com/questions/114151/how-do-i-reference-in-appendix-a-theorem-given-in-the-body
\usepackage{thmtools, thm-restate}

%% Theoretical computer science
\usepackage{stmaryrd}
\usepackage{mathtools}  % For "::=" ( \Coloneqq )

%% Font
% \usepackage[euler-digits,euler-hat-accent]{eulervm}

\usepackage{pifont}
\usepackage{textgreek}



\usepackage{savesym}
\savesymbol{checkmark}
\usepackage{dingbat}

\usepackage{comment}

\usepackage{longtable}

\usepackage{tikz-cd}
\usetikzlibrary{babel}


% Code highlighting
\usepackage{listings}

\lstset{%
  % backgroundcolor=\color{white},
  basicstyle=\small\ttfamily,
  keywordstyle=\sffamily\bfseries,
  captionpos=none,
  columns=flexible,
  keepspaces=true,
  showspaces=false,               % show spaces adding particular underscores
  showstringspaces=false,         % underline spaces within strings
  showtabs=false,                 % show tabs within strings adding particular underscores
  breaklines=true,                % sets automatic line breaking
  breakatwhitespace=true,         % sets if automatic breaks should only happen at whitespace
  escapeinside={(*}{*)},
  literate={lam}{{$\lambda$}}1 {->}{{$\rightarrow$}}1 {Top}{{$\top$}}1 {o+o}{{$\oplus$}}1 {=>}{{$\Rightarrow$}}1 {/\\}{{$\Lambda$}}1 {forall}{{$\forall$}}1,
  tabsize=2,
  commentstyle=\color{purple}\ttfamily,
  stringstyle=\color{red}\ttfamily,
  sensitive=false
}

\lstdefinelanguage{sedel}{
  keywords={main, Int, String, Bool, self, Void, Class, All, extends, this, trait, inherits, super, type, Trait, override, new, if, then, else, let, in, letrec},
  identifierstyle=\color{black},
  morecomment=[l]{--},
  morecomment=[l]{//},
  morestring=[b]",
  xleftmargin  = 3mm,
  morestring=[b]'
}


\lstdefinelanguage{simple}{
  keywords={letrec, in, let},
  identifierstyle=\color{black},
  xleftmargin  = 3mm,
}


\lstdefinelanguage{gbeta}{%
  language     = java,
  morekeywords = {virtual,refine},
  xleftmargin  = 3mm
}


\lstdefinelanguage{JavaScript}{
  keywords={function, const, extends, super, class, export, boolean, throw, implements, import, this, typeof, new, true, false, catch, function, return, null, catch, switch, var, if, in, while, do, else, case, break},
  identifierstyle=\color{black},
  comment=[l]{//},
  morecomment=[s]{/*}{*/},
  morestring=[b]',
  xleftmargin  = 3mm,
  morestring=[b]"
}


\lstset{language=sedel}

\theoremstyle{plain}
\newtheorem{theorem}{Theorem}[chapter]
\newtheorem{lemma}[theorem]{Lemma}
\newtheorem{corollary}[theorem]{Corollary}
\newtheorem{proposition}[theorem]{Proposition}
\newtheorem{conjecture}[theorem]{Conjecture}
\theoremstyle{definition}
\newtheorem{definition}{Definition}
\newtheorem{example}{Example}
\newtheorem{observation}{Observation}
\theoremstyle{remark}
\newtheorem*{remark}{Remark}


%%%%%%%%%%%%%%%%%%%%%%%%%%%%%%%%%%%%%%%%%%%%%%%%%%%%%%%%%%%%%%%%%%%%%%%%
% Ott includes
%%%%%%%%%%%%%%%%%%%%%%%%%%%%%%%%%%%%%%%%%%%%%%%%%%%%%%%%%%%%%%%%%%%%%%%%

\usepackage{ottalt}
\inputott{Gen/ott-rules}
\renewcommand\ottaltinferrule[4]{
  \inferrule*[narrower=0.6,lab=#1,#2]
    {#3}
    {#4}
}
%%----------------------------------------------------------------------------%%

%%----------------------------------------------------------------------------%%
%%    Environment for Title Page                                              %%
%%----------------------------------------------------------------------------%%

\makeatletter
\renewcommand{\maketitle}
{%
  \begin{titlepage}
    \renewcommand{\baselinestretch}{1}
    \begin{center}
      \vspace*{\stretch{3}}
      {\LARGE\@title\par}
      \vspace*{1cm}
      {\large\textit{by}\par}
      \vspace*{1cm}
      {\Large\@author\par}
      {\Large(谢宁宁)\par}
      \vspace*{\stretch{5}}
      {{\includegraphics[width=30mm]{figures/hku}} \par}
      {\hbox{}\par}
      \vspace*{\stretch{5}}
      {
      	% \textsl{(Temporary Binding for Examination Purposes)} \\
      	% \vspace*{3cm}
      	{\normalsize
      	A thesis submitted in partial fulfillment of the requirements for \\
        the degree of Doctor of Philosophy \\
        at The University of Hong Kong \\
        \par
        }
      }
      \vspace*{\stretch{1}}
      {\large\@date\par}
      \vspace*{\stretch{1}}
    \end{center}
  \end{titlepage}
} \makeatother




%%----------------------------------------------------------------------------%%
%%    Environment for Abstract                                                %%
%%----------------------------------------------------------------------------%%

\makeatletter
\newenvironment{abstract}%
{
  \if@twocolumn
    \@restonecoltrue\onecolumn
  \else
    \@restonecolfalse\newpage
  \fi
  \pagestyle{empty}%
  \setcounter{page}\@ne
  \mbox{}
  \newpage
  \pagestyle{empty}%
  \setcounter{page}\@ne
  \cleardoublepage
  \begin{center}
  \large
  Abstract of thesis entitled            \par
  \textbf{``\@title''}                   \par
                                         \vspace*{0.2in}
  Submitted by                           \par
  \@author                               \par
                                         \vspace*{0.2in}
  for the degree of Doctor of Philosophy \\
  at The University of Hong Kong         \\
  in \@date
  \end{center}
  \vspace*{0.5in}
}%
{
  \if@restonecol\twocolumn \else \newpage \fi
  \if@twoside\else
  \setcounter{page}\@ne
  \fi
}
\makeatother

%%----------------------------------------------------------------------------%%
%%    Environment for Declaration                                             %%
%%----------------------------------------------------------------------------%%


\makeatletter
\newcommand{\makedeclaration}
{
\chapter*{Declaration}
\addcontentsline{toc}{chapter}{Declaration}
\noindent I declare that this thesis represents my own work, except
where due acknowledgment is made, and that it has not been
previously included in a thesis, dissertation or report submitted
to this University or to any other institution for a degree,
diploma or other qualifications.
\vspace*{1.5in}

\noindent%
\begin{tabular}{@{}l@{}}
\dotfill \\
\@author\hspace*{3cm}\\
\@date\\
\end{tabular}
}
\makeatother

%%----------------------------------------------------------------------------%%
%%    Environment for Acknowledgments                                         %%
%%----------------------------------------------------------------------------%%

\makeatletter
\newcommand{\makeAck}
{
\chapter*{Acknowledgments}
\addcontentsline{toc}{chapter}{Acknowledgments}
I am incredibly fortunate to have many people to thank.

First and foremost, I would like to thank my advisor Bruno C. d. S. Oliveira,
who guided me into the land of functional programming. He was always very
patient with me, no matter what crazy research ideas I would like to try. Every
meeting I had with Bruno was fruitful -- he could consistently provide insightful
suggestions and was also willing to dive deeply into technical discussions. It
is his mentorship and encouragement which made this thesis possible.

I thank my thesis committee, Zhiyi Huang and Giulio Chiribella, for their careful
reading of the draft, and especially Simon Peyton Jones for his 18-page long
review and various constructive feedback -- the thesis has improved a lot
because of his suggestions.

I would like to thank Tom Schrijvers, who hosted me during my visit to KU
Leuven. Tom is insightful and has the fantastic ability to explain complicated
things in an abstract and easy-to-follow manner. I am also very grateful for his
long-lasting support with my research career.

Richard Eisenberg introduced me to the Haskell community. I greatly enjoyed my
summer stay in Philadelphia when hacking GHC. I always try to learn from him how
to communicate research ideas effectively and how to think practically.

Great thanks to Daan Leijen, who hosted me during my internship at Microsoft
Research. Daan has lots of exciting ideas, and his research vision on Koka has
always inspired me. Our collaboration has been so fruitful, and I expect much
more to come. I also thank Madan Musuvathi and the whole RiSE team for their
generous support of this internship which happened during the pandemic.

I am grateful to the wonderful programming language friends I have met all over
the world. I had lots of fun with the PLGroup at HKU: Xuan Bi, Xuejing Huang,
Huang Li, Baber Rehman, Tomas Tauber, Yanlin Wang, Yanpeng Yang, Weixin Zhang,
Jinxu Zhao, Yaoda Zhou. Especially Xuan Bi, who has co-authored several papers
with me, and his passion for Haskell was a big inspiration for me during my
early Ph.D. years. I also thank Tom's excellent students Gert-Jan Bottu, George
Karachalias, Steven Keuchel, Klara Mardirosian, Koen Pauwels, Ruben Pieters, Amr
Saleh, César Santos, and Alexander Vandenbroucke, who made my stay at KU Leuven
comfortable and fun. Yao Li and Yishuai Li have helped me a lot during my stay
in Philadelphia. My other great collaborators include Jonathan Brachthäuser,
Daniel Hillerström, Leonardo de Moura,  Alex Reinking, and Philipp Schuster.

Last but not least, I would like to thank my parents, my brother, and Xujie Si
for their unconditional love and support. I made it through thanks to all of
you.
}
\makeatother


\KOMAoptions{listof=totoc}
% \setkomafont{sectioning}{\normalfont\bfseries}
